\clearpage
\section{Poiseuille Flow Between Concentric Cylinders}
\begin{colbox}
    Assume that a fluid flows between two infinite concentric cylinders with radii \(\ins{r_2 > r_1}\), which are at rest. The flow is parallel to the symmetry axis \(\vu{z}\). The fluid is Newtonian with viscosity \(\mu\) (satisfying the Navier-Sokes equations). Assume there are no external forces and that the pressure gradient is constant and also directed along the symmetry axis: \(\ins{\del{z}p=\text{const.}}\)
    \begin{enumerate}
        \item Write the full equation of motion in cylindrical coordinates.
        \item Use the symmetries of the problem to sinmplife the equation. What is the form of the nonlinear term \(\ins{\vb{v}\cdot\nabla}\vb{v}\)?
        \item Write the boundary conditions of the problem and solve the resulting system to find the velocity profile.
        \item Define and compute the average velocity through a cross-sectional area, \(\overline{U}\). Use the average velocity to define the Reynolds number Re and the momentum flux through a cross-section. Explain the meaning of these two quantities.
        \item Compute the friction force (per unit area) acting on the cylinder walls and express it in dimensionless form using the momentum flux you defined. What is the dependence of the result on Re? Explain this dependence.
        \item Use your solutions from parts 3-5 to find the flow profile and the friction force in a hollow cylinder.
    \end{enumerate}
\end{colbox}
\subsection{Equations of motion}
Since the fluid is Newtonian, the NS equation relevant to us is
\begin{equation}
    \rho\ins{\del{t}+\ins{\vb{v}\cdot\nabla}}\vb{v} = -\grad{p}+\mu\laplacian{\vb{v}}+\rho\vb{f}
\end{equation}
inserting the definitions of the diferential operators in cylindrical
\begin{equation}
    \rho\ins{\del{t}+\ins{\vb{v}\cdot\ins{\pdv{r}\vu{r}+\frac{1}{r}\pdv{\theta}\vu{\theta}+\pdv{z}\vu{z}}}}\vb{v} = -\cdot\ins{\pdv{r}\vu{r}+\frac{1}{r}\pdv{\theta}\vu{\theta}+\pdv{z}\vu{z}}p+\mu\ins{\frac{1}{r}\pdv{r}\ins{r\pdv{r}}+\frac{1}{r^2}\pdv[2]{\theta}+\pdv[2]{z}}v+\rho\vb{f}
\end{equation}
which is awful.
\subsection{simplification}
Since we assume the flow is parallel to \(\vu{z}\), there is no external force, and the pressure gradient is constant, we can simplify to
\begin{equation}
    \rho\ins{\del{t}+v_z\del{z}}\vb{v} = -\del{z}p+\mu\laplacian{\vb{v}}
\end{equation}
To simplify the laplacian we can note that the velocity clearly can't have dependence on the radial direction or the z position, from symmetry, therefore
\begin{equation}
    \rho\ins{\del{t}}\vb{v} = -\del{z}p+\mu\ins{\frac{1}{r}\del{r}\ins{r\del{r}}}\vb{v}
\end{equation}
From symmetry and conservation laws, the flow must remain constant in time, therefore
\begin{equation}
    0 = -\del{z}p+\mu\ins{\frac{1}{r}\del{r}\ins{r\del{r}}}\vb{v}
\end{equation}
rearranging
\begin{equation}
    \del{z}p = \mu\ins{\frac{1}{r}\del{r}\ins{r\del{r}}}\vb{v}
\end{equation}
\subsection{Boundary Conditions}
Since both cylinders are stationary we must demand that 
\begin{align}
    v(r_1)=v(r_2)=0
\end{align}
We'll rearrange the equation to be for the velocity
\begin{align}
     \ins{\frac{1}{r}\del{r}\ins{r\del{r}}}\vb{v}=\frac{\del{z}p}{\mu} \\
     \ins{\del{r}\ins{r\del{r}}}\vb{v}=\frac{r\del{z}p}{\mu} \\
     r\del{r}v = \frac{\del{z}p}{2\mu}r^2+B \\
     v = \frac{\del{z}p}{4\mu}r^2+B\ln(r)+C
\end{align}
using the boundary conditions
\begin{align}
    \frac{\del{z}p}{4\mu}r_1^2+B\ln(r_1)+C &= \frac{\del{z}p}{4\mu}r_2^2+B\ln(r_2)+C = 0 \\
    B\ln(r_1)-B\ln(r_2) &= \frac{\del{z}p}{4\mu}r_2^2-\frac{\del{z}p}{4\mu}r_1^2 \\
    B\ln(\frac{r_1}{r_2})&=\frac{\del{z}p}{4\mu}\ins{r_2^2-r_1^2} \\
    B &= \frac{\del{z}p\ins{r_2^2-r_1^2}}{4\mu\ln(\frac{r_1}{r_2})}
\end{align}
inserting back into one of the equations
\begin{align}
    \frac{\del{z}p}{4\mu}r_1^2+\frac{\del{z}p\ins{r_2^2-r_1^2}}{4\mu\ln(\frac{r_1}{r_2})}\ln(r_1)+C = 0 \\
    \frac{\del{z}p}{4\mu}\ins{r_1^2+\frac{\ins{r_2^2-r_1^2}}{\ln(\frac{r_1}{r_2})}\ln(r_1)}+C = 0 \\
    C = -\frac{\del{z}p}{4\mu}\ins{r_1^2+\frac{\ins{r_2^2-r_1^2}}{\ln(\frac{r_1}{r_2})}\ln(r_1)}
\end{align}
\subsection{average velocity}
